\documentclass[UTF8,zihao=-4]{ctexart}
\usepackage[a4paper,margin=2.2cm]{geometry}
\usepackage{fontspec}
\usepackage{tabularx}
\usepackage{array}
\usepackage{ragged2e}
\usepackage{fancyhdr}
\usepackage{hyperref}
\setCJKmainfont{Noto Serif SC}
\setCJKsansfont{Noto Sans SC}
\setmainfont{Noto Serif SC}
\setlength{\parindent}{0pt}
\setlength{\parskip}{0.45em}
\pagestyle{fancy}
\fancyhf{}
\fancyhead[L]{商务智能}
\fancyfoot[C]{\thepage}
\hypersetup{colorlinks=true,linkcolor=black,urlcolor=blue}
\begin{document}
\title{9、数据刷新的四种方法}
\author{}
\date{}
\maketitle

时间戳、DELTA 文件、建立映象文件、日志文件\par

数据刷新主要有四种方法：时间戳、DELTA 文件、建立映象文件和日志文件。时间戳依赖记录中的修改时间；DELTA 文件记录应用修改操作；映象文件通过两次快照比较差异；日志文件利用 OLTP 数据库日志进行增量刷新。四种方法可以在同一数据仓库系统中结合使用，以适应不同数据源。\par

\end{document}
